\section{The present analytic boundary}
\label{sec:tpc45-literature}

The congruence and determinant reformulations above place the remaining
TPC term near several established dispersion and
Kloosterman--fraction estimates.  The resemblance is exact at the level
of the equation, but not at the level of the coefficient class or of the
available exponent range.  We record the distinction explicitly.  In
this section \(Q_{\mathrm{mod}}\) denotes a modulus scale; it is not the
row-length parameter used elsewhere in the paper.

\subsection{Unbalanced-convolution dispersion}

Fouvry and Radziwi{\l}{\l} prove distribution estimates for a convolution
\(
  \sum_{mn\equiv a\, (\bmod q)}\alpha_m\beta_n
\)
averaged over \(q\asymp Q_{\mathrm{mod}}\), with divisor-bounded
coefficients and a Siegel--Walfisz hypothesis on the short sequence
\cite[Corollary~1.1]{FR2022}.  With \(MN\asymp X\), their first range is
\begin{equation}
 \exp\!\bigl((\log X)^\epsilon\bigr)
 \le N
 \le Q_{\mathrm{mod}}^{-11/12}X^{17/36-\epsilon}.
 \label{eq:tpc45-fr-first-range}
\end{equation}
Their two other stated ranges have, respectively,
\begin{equation}
 \begin{aligned}
  N&\le X^{7/90-\epsilon},
  &Q_{\mathrm{mod}}&\le X^{53/105-\epsilon},\\
  N&\le X^{101/630-\epsilon},
  &Q_{\mathrm{mod}}&\le X^{53/105-\epsilon},
 \end{aligned}
 \label{eq:tpc45-fr-other-ranges}
\end{equation}
with the further hypotheses specified there.  Their prime-modulus
variant removes the Siegel--Walfisz input but retains the power range in
\eqref{eq:tpc45-fr-first-range} \cite[Theorem~1.2]{FR2022}.

The algebraic match with our incidence relation is literal:
\[
 p\mid(mj+h)
 \quad\Longleftrightarrow\quad
 mj\equiv-h\pmod p.
\]
The parameter match fails, however.  At the first large-prime block
\(p\asymp Q_{\mathrm{mod}}=X^{1/2}\),
\eqref{eq:tpc45-fr-first-range} permits only
\[
 N\le X^{17/36-11/24-\epsilon}
   =X^{1/72-\epsilon},
\]
whereas the TPC orbit variable has length
\[
 J=X^{133/400+o(1)}.
\]
The exponent deficit is
\begin{equation}
 \frac{133}{400}-\frac1{72}
 =\frac{1147}{3600}.
 \label{eq:tpc45-fr-exponent-deficit}
\end{equation}
The alternatives in \eqref{eq:tpc45-fr-other-ranges} are also shorter than
\(J\).  Moreover, successive physical prime restrictions produce
coefficient fields depending on the preceding restrictions, rather than
a single fixed product \(\alpha_m\beta_j\).  Thus the theorem has the
right congruence but neither the required length nor the required
restriction-stable coefficient interface.

Wright's partially fixed-modulus Kloosterman estimate improves the first
unbalanced range to
\begin{equation}
 N\le Q_{\mathrm{mod}}^{-33/28}X^{17/28-\epsilon}
 \label{eq:tpc45-wright-first-range}
\end{equation}
and improves the accompanying modulus cap from \(53/105\) to \(45/89\)
in the other two cases \cite[Corollary~2.2]{Wright2026}.  At
\(Q_{\mathrm{mod}}=X^{1/2}\), \eqref{eq:tpc45-wright-first-range} becomes
\[
 N\le X^{17/28-33/56-\epsilon}
   =X^{1/56-\epsilon}.
\]
Consequently the gap to the TPC length is still
\begin{equation}
 \frac{133}{400}-\frac1{56}
 =\frac{881}{2800}.
 \label{eq:tpc45-wright-exponent-deficit}
\end{equation}
The advance is genuine, but it does not enter the balanced TPC regime.

\subsection{The fixed-determinant coefficient-slot barrier}

Bettin and Chandee estimate trilinear forms with Kloosterman fractions
for three arbitrary coefficient sequences \cite[Theorem~1]{BC2018}.
Their fixed-determinant consequence concerns
\begin{equation}
 m_1n_2-m_2n_1=\Delta
 \label{eq:tpc45-bc-fixed-determinant}
\end{equation}
and places arbitrary arithmetic sequences on \(n_1,n_2\), while the
weights on \(m_1,m_2\) are smooth dyadic functions with controlled
derivatives \cite[Corollary~1]{BC2018}.  Hence this corollary has two,
not three, arbitrary arithmetic coefficient slots.

Our cofactor equation
\begin{equation}
 pt-mj=h
 \label{eq:tpc45-literature-determinant}
\end{equation}
has the same determinant shape.  Its actual coefficient system contains
the prime restriction in \(p\), the M\"obius weight \(\mu(t)\), and the
source weight \(\mu(d_\rho)\) attached to
\(m_\rho=\ell_\rho d_\rho\), together with the inherited masks in \(m_\rho\).
Only the orbit variable \(j\) is naturally smooth.  Assigning two of
\(p,t,m\) to the arbitrary slots in \eqref{eq:tpc45-bc-fixed-determinant}
therefore leaves the third nonsmooth arithmetic sequence in a slot for
which the corollary requires a smooth weight.  The three-sequence
trilinear theorem does not by itself restore this lost slot in the
fixed-\(h\) determinant corollary.  This is the precise
\emph{third-coefficient barrier}; it is not merely a difference of
notation.

\begin{remark}[A diagnostic exponent ledger, not a theorem]
\label{rem:tpc45-bc-ledger-not-theorem}
It is informative to make a deliberately counterfactual calculation.
Suppose that two TPC variables could be replaced by admissible smooth
weights and that the remaining arbitrary variables had lengths
\[
 P=X^{267/400},\qquad T=X^{133/400},\qquad PT=X.
\]
Substituting these lengths into the error term of
\cite[Corollary~1]{BC2018} produces the formal exponent
\[
 \frac12+\frac7{20}
 +\frac14\frac{267}{400}
 =1+\frac{27}{1600}.
\]
A formal comparison with an \(X^{1+1/400+o(1)}\) endpoint budget then
leaves \(23/1600\) in the exponent.  This computation is
\emph{not an application} of the Bettin--Chandee corollary: the
counterfactual smoothing discards an actual arithmetic weight, and a
scalar determinant-count error is not automatically the vector
Gram/Bessel energy required here.  It is neither a lower bound nor an
impossibility theorem.  It records only that the most naive
two-arbitrary-slot substitution would not close the desired exponent
ledger.
\end{remark}

\subsection{A withdrawn bilinear claim}

The first version of the 2026 preprint of Dong, Robles and Zeindler
announced an improved bilinear Kloosterman--fraction bound and a saving
of \(1/12\) over the trivial estimate in the balanced case
\cite[Theorems~1.4--1.6]{DRZ2026}.  The authors subsequently withdrew
the paper.  The withdrawal notice states that a missing factor in one
equation changes \(L^5\) to \(L^7\), so the argument no longer yields
the claimed improved bound \cite{DRZ2026}.  We therefore do not use the
announced \(1/12\) saving as a theorem or as an input anywhere in this
paper.

There is a second, logically independent boundary.  Even if the
announced bilinear estimate were available, it concerns a form with two
arbitrary sequences and a Kloosterman phase.  No fixed-\(h\)
determinant corollary with the three arithmetic weights in
\eqref{eq:tpc45-literature-determinant} is supplied.  Thus that announced
estimate would improve a possible analytic subroutine, but would not by
itself cross the third-coefficient interface.

\subsection{Why the Bourgain--Sarnak--Ziegler criterion is not a black box}

The Bourgain--Sarnak--Ziegler criterion gives a powerful route from
prime-dilate decorrelation to multiplicative orthogonality.  In one
form, if \(|F|\le1\), \(|\nu|\le1\) is multiplicative, and
\[
 \left|\sum_{m\le M}F(p_1m)\overline{F(p_2m)}\right|
 \le \tau M
\]
uniformly for the relevant distinct primes \(p_1,p_2\), then
\[
 \left|\sum_{n\le N}\nu(n)F(n)\right|
 \le 2\sqrt{\tau\log(1/\tau)}\,N
\]
for sufficiently large parameters \cite[Theorem~2]{BSZ2013}.

This criterion identifies the right qualitative phenomenon, but it is
not directly applicable to the TPC coefficient.  Here the coefficient
is sparse and depends on the scale \(X\); its normalization is governed
by a projective energy rather than by the ambient interval length.
More importantly, after a physical restriction at one prime, the
coefficient used at the next prime is a new restricted field, so the
needed correlations involve restriction-dependent pairs rather than
the dilates of one fixed bounded function \(F\).  Dividing by an
\(\ell^\infty\) norm and applying the criterion would return an
ambient-length estimate, not the endpoint Gram-energy bound.  A usable
BSZ-type reduction would therefore first require a new
restriction-stable, energy-normalized prime-dilate theorem; the cited
criterion does not supply it.

\subsection{Exact frontier statement}

The preceding comparisons do not show that
\eqref{eq:tpc45-literature-determinant} is intractable.  They show only
that the following interface is absent from the cited black boxes:
\begin{enumerate}[label=\textup{(\roman*)},leftmargin=2.4em]
 \item the shift \(h\) is fixed rather than averaged;
 \item the prime, cofactor-M\"obius, and source-M\"obius/mask weights
       remain simultaneously present;
 \item the estimate is stable under the projective and physical
       restrictions inherited from the TPC packet; and
 \item the output controls a vector-valued Bessel/Gram synthesis norm,
       not only a scalar congruence discrepancy.
\end{enumerate}
Any future positive input must address these four points explicitly.
An estimate averaged over \(h\), or one obtained after replacing a
nonsmooth coefficient by a smooth weight, may be valuable in its own
right but does not close the fixed physical TPC gate.
